%------------------------------------------------------------------------------%
\documentclass[fleqn,utf8,aspectratio=169,12pt]{beamer}

\usepackage{gaal}

\title{Projeção Ortogonal}

%------------------------------------------------------------------------------%
\begin{document}

\addtitle

%------------------------------------------------------------------------------%
\section{Decomposição de Vetores}

%------------------------------------------------------------------------------%
\begin{frame}{Decomposição de Vetores}
  Pode ser conveniente decompor o vetor em uma soma de vetores

  \pause
  \bigskip
  Por exemplo,

  \pause
  Podemos decompor qualquer vetor em vetores alinhados com os eixos do $\R^n$
\end{frame}

%------------------------------------------------------------------------------%
\framedgraphic*{Decomposição de Vetores}{E-07-Projecao-decomposicao-1}{}
\framedgraphic*{Decomposição de Vetores}{E-07-Projecao-decomposicao-2}{}
\framedgraphic*{Decomposição de Vetores}{E-07-Projecao-decomposicao-3}{}

%------------------------------------------------------------------------------%
\begin{frame}{Decomposição de Vetores}
  Muitas vezes precisamos decompor um vetor na direção de outo

  \pause
  \bigskip
  Para isso precisamos da \emph{projeção ortogonal} de um vetor em outro
\end{frame}

%------------------------------------------------------------------------------%
\framedgraphic*{Decomposição Ortogonal}{E-07-Projecao-decom-ortogonal-1}{}
\framedgraphic*{Decomposição Ortogonal}{E-07-Projecao-decom-ortogonal-2}{}
\framedgraphic*{Decomposição Ortogonal}{E-07-Projecao-decom-ortogonal-3}{}
\framedgraphic*{Decomposição Ortogonal}{E-07-Projecao-decom-ortogonal-4}{}

%------------------------------------------------------------------------------%
\section{Projeção Ortogonal}

%------------------------------------------------------------------------------%
\begin{frame}{Projeção ortogonal}
  Considere os vetores $\vec{w}$ e $\vec{v}$

  \pause
  Queremos determinar a componente de $\vec{w}$ na direção de $\vec{v}$

  \pause
  A \emph{projeção ortogonal} de $\vec{w}$ sobre $\vec{v}$

  \pause
  é o vetor paralelo a $\vec{v}$

  \pause
  obtido ao projetar ortogonalmente $\vec{w}$

  \pause
  sobre a reta determinada por $\vec{v}$
\end{frame}

%------------------------------------------------------------------------------%
\begin{frame}{Projeção Vetorial}
\begin{columns}
\column{0.5\textwidth}
  \onslide<+->{Se $\theta$ é o ângulo entre $\vec{w}$ e $\vec{v}$}

  \onslide<+->{A componente de $\vec{w}$ na direção de $\vec{v}$ \\}
  \onslide<+->{possui comprimento
  \[
    \norm{\vec{w}}\cos(\theta)
  \]}
  \onslide<+->{Do produto escalar temos}
  \begin{align*}
    \onslide<+->{\vec{w}\cdot\vec{v} & = \norm{\vec{w}} \norm{\vec{v}} \cos(\theta) \\}
    \onslide<+->{\norm{\vec{w}} \cos(\theta) & = \frac{\vec{w}\cdot\vec{v}}{\norm{\vec{v}}}}
  \end{align*}
\column{0.4\textwidth}
  \includegraphics{E-07-Projecao-pytagoras}
\end{columns}
\end{frame}

%------------------------------------------------------------------------------%
\begin{frame}{Projeção Vetorial}
  Queremos agora um vetor na direção de $\vec{v}$,
  \pause
  com comprimento
  \(
    \frac{\vec{w}\cdot\vec{v}}{\norm{\vec{v}}}
  \)

  \pause
  O vetor unitário na direção de $\vec{v}$ é
  \(
    \frac{\vec{v}}{\norm{\vec{v}}}
  \)

  \pause
  Portanto
  \[
    \pause
    \projuv{w}{v}
    \pause
    \;=\; \left( \frac{\vec{w}\cdot\vec{v}}{\norm{\vec{v}}} \right)
          \frac{\vec{v}}{\norm{\vec{v}}}
    \pause
    \;=\; \frac{\vec{w}\cdot\vec{v}}{\norm{\vec{v}}^2} \, \vec{v}
    \pause
    \;=\; \frac{\vec{w}\cdot\vec{v}}{\vec{v}\cdot\vec{v}} \, \vec{v}
  \]

  \pause
  desde que
  \(
    \vec{v} \neq \vec{0}
  \)
\end{frame}

%------------------------------------------------------------------------------%
\begin{frame}{Projeção Vetorial}
  Projeção do vetor \({\color{blue}w}\)
  na direção do vetor \({\color{magenta}v} \neq \vec{0} \)
  \Large
  \[
    \projuv{{\color{blue}w}}{{\color{magenta}v}}
    \;=\; \frac{\vec{{\color{blue}w}}\cdot\vec{{\color{magenta}v}}}
               {\vec{{\color{magenta}v}}\cdot\vec{{\color{magenta}v}}}
               \, \vec{{\color{magenta}v}}
  \]
\end{frame}

%------------------------------------------------------------------------------%
%------------------------------------------------------------------------------%
\begin{question}
  Dados
  \[
    \vec{w} = (3, 4)
    \qquad
    \vec{v} = (1, 0)
  \]
  Calcule a projeção de $\vec{w}$ sobre $\vec{v}$
\end{question}

%------------------------------------------------------------------------------%
\begin{resolution}{Solução}
  Queremos a projeção de $\vec{w}$ em $\vec{v}$
  \pause
  \[
    \projuv{w}{v}
    \;=\; \frac{\vec{w}\cdot\vec{v}}{\vec{v}\cdot\vec{v}} \, \vec{v}
  \]
\end{resolution}

%------------------------------------------------------------------------------%
\begin{resolution}{Solução}
\begin{columns}[t]
\column{0.5\textwidth}
  \begin{align*}
    \onslide<+->{\vec{w} \cdot \vec{v}}
    \onslide<+->{& = (3, 4) \cdot (1, 0) \\}
    \onslide<+->{& = 3 \times 1 + 4 \times 0 \\}
    \onslide<+->{& = 3}
  \end{align*}
\column{0.5\textwidth}
  \begin{align*}
    \onslide<+->{\vec{v}\cdot\vec{v}}
    \onslide<+->{& = (1, 0) \cdot (1, 0) \\}
    \onslide<+->{& = 1 \times 1 + 0 \times 0 \\}
    \onslide<+->{& = 1}
  \end{align*}
\end{columns}
  \vspace{-\baselineskip}
  \[
    \onslide<+->{\projuv{w}{v}}
    \onslide<+->{\;=\; \frac{\vec{w}\cdot\vec{v}}{\vec{v}\cdot\vec{v}} \, \vec{v}}
    \onslide<+->{\;=\; \frac{3}{1} \vetor{1}{0}}
    \onslide<+->{\;=\; \vetor{3}{0}}
  \]
\end{resolution}

%------------------------------------------------------------------------------%
\endinput

%------------------------------------------------------------------------------%
\begin{question}
  Dados
  \[
    \vec{w} = (3, 4)
    \qquad
    \vec{v} = (2, 1)
  \]
  Calcule a projeção de $\vec{w}$ sobre $\vec{v}$
\end{question}

%------------------------------------------------------------------------------%
\begin{resolution}{Solução}
  Queremos a projeção de $\vec{w}$ em $\vec{v}$
  \pause
  \[
    \projuv{w}{v}
    \;=\; \frac{\vec{w}\cdot\vec{v}}{\vec{v}\cdot\vec{v}} \, \vec{v}
  \]
\end{resolution}

%------------------------------------------------------------------------------%
\begin{resolution}{Solução}
\begin{columns}[t]
\column{0.5\textwidth}
  \begin{align*}
    \onslide<+->{\vec{w} \cdot \vec{v}}
    \onslide<+->{& = (3, 4) \cdot (2, 1) \\}
    \onslide<+->{& = 3 \times 2 + 4 \times 1 \\}
    \onslide<+->{& = 6 + 4 \\}
    \onslide<+->{& = 10}
  \end{align*}
\column{0.5\textwidth}
  \begin{align*}
    \onslide<+->{\vec{v}\cdot\vec{v}}
    \onslide<+->{& = (2, 1) \cdot (2, 1) \\}
    \onslide<+->{& = 2^2 + 1^2 \\}
    \onslide<+->{& = 4 + 1 \\}
    \onslide<+->{& = 5}
  \end{align*}
\end{columns}
  \vspace{-\baselineskip}
  \[
    \onslide<+->{\projuv{w}{v}}
    \onslide<+->{\;=\; \frac{\vec{w}\cdot\vec{v}}{\vec{v}\cdot\vec{v}} \, \vec{v}}
    \onslide<+->{\;=\; \frac{10}{5} \vetor{2}{1}}
    \onslide<+->{\;=\; 2 \vetor{2}{1}}
    \onslide<+->{\;=\; \vetor{4}{2}}
  \]
\end{resolution}

%------------------------------------------------------------------------------%
\endinput

%------------------------------------------------------------------------------%
\begin{question}
  Dados
  \[
    \vec{w} = (3, 4)
    \qquad
    \vec{v} = (2, 1)
  \]
  Calcule a projeção de $\vec{v}$ sobre $\vec{w}$

  \pause
  \bigskip
  \emph{Atenção à ordem dos vetores}
\end{question}

%------------------------------------------------------------------------------%
\begin{resolution}{Solução}
  Queremos a projeção de $\vec{v}$ em $\vec{w}$
  \pause
  \[
    \projuv{v}{w}
    \;=\; \frac{\vec{v}\cdot\vec{w}}{\vec{w}\cdot\vec{w}} \, \vec{w}
  \]
\end{resolution}

%------------------------------------------------------------------------------%
\begin{resolution}{Solução}
\begin{columns}[t]
\column{0.5\textwidth}
  \begin{align*}
    \onslide<+->{\vec{v} \cdot \vec{w}}
    \onslide<+->{= \vec{w} \cdot \vec{v}}
    \onslide<+->{= 10}
  \end{align*}
\column{0.5\textwidth}
  \begin{align*}
    \onslide<+->{\vec{w}\cdot\vec{w}}
    \onslide<+->{& = (3, 4) \cdot (3, 4) \\}
    \onslide<+->{& = 3^2 + 4^2 \\}
    \onslide<+->{& = 9 + 16 \\}
    \onslide<+->{& = 25}
  \end{align*}
\end{columns}
  \[
    \onslide<+->{\projuv{v}{w}}
    \onslide<+->{\;=\; \frac{\vec{v}\cdot\vec{w}}{\vec{w}\cdot\vec{w}} \, \vec{w}}
    \onslide<+->{\;=\; \frac{10}{25} \vetor{3}{4}}
    \onslide<+->{\;=\; \frac{2}{5} \vetor{3}{4}}
    \onslide<+->{\;=\; \Vetor{\frac{6}{5}}{\frac{8}{5}}}
  \]
\end{resolution}

%------------------------------------------------------------------------------%
\endinput


%------------------------------------------------------------------------------%
\begin{frame}{Decomposição Ortogonal}
  \onslide<+->{Um vetor $\vec{w}$ pode ser decomposto em duas componentes}
  \[
    \onslide<+->{\vec{w}}
    \onslide<+->{= \vec{w}_{\parallel} + \vec{w}_{\perp}}
  \]
  \onslide<+->{onde}
  \begin{align*}
    \onslide<+->{\vec{w}_{\parallel\;} & = \projuv{w}{v}}
    & &
    \onslide<7->{\vec{w}_{\parallel\;\,} \text{ é paralelo a } \vec{v}}
    \\
    \onslide<+->{\vec{w}_{\perp}     & = \vec{w}-\vec{w}_{\parallel}}
    & &
    \onslide<8->{\vec{w}_{\perp} \text{ é ortogonal a } \vec{v}}
  \end{align*}
\end{frame}

%------------------------------------------------------------------------------%
%------------------------------------------------------------------------------%
\begin{question}
  Dados
  \[
    \vec{w} = \vetor{5}{2}
    \qquad
    \vec{v} = \vetor{1}{2}
  \]
  Determine a componente de $\vec{w}$ perpendicular a $\vec{v}$
\end{question}

%------------------------------------------------------------------------------%
\begin{resolution}{Solução}
  Queremos a \emph{componente perpendicular} de $\vec{w}$ em $\vec{v}$
  \pause
  \[
    \vec{w}_\perp
    \;=\; \vec{w} - \projuv{w}{v}
    \;=\; \vec{w} - \frac{\vec{w}\cdot\vec{v}}{\vec{v}\cdot\vec{v}} \, \vec{v}
  \]
\end{resolution}

%------------------------------------------------------------------------------%
\begin{resolution}{Projeção Ortogonal}
\begin{columns}[t]
\column{0.5\textwidth}
  \begin{align*}
    \onslide<+->{\vec{w} \cdot \vec{v}}
    \onslide<+->{& = \vetor{5}{2} \cdot \vetor{1}{2} \\}
    \onslide<+->{& = 5 \times 1 + 2 \times 2\\}
    \onslide<+->{& = 5 + 4 \\}
    \onslide<+->{& = 9}
  \end{align*}
\column{0.5\textwidth}
  \begin{align*}
    \onslide<+->{\vec{v}\cdot\vec{v}}
    \onslide<+->{& = \vetor{1}{2} \cdot \vetor{1}{2} \\}
    \onslide<+->{& = 1^2 + 2^2 \\}
    \onslide<+->{& = 1 + 4 \\}
    \onslide<+->{& = 5 }
  \end{align*}
\end{columns}
  \vspace{-\baselineskip}
  \[
    \onslide<+->{\projuv{w}{v}}
    \onslide<+->{\;=\; \frac{\vec{w}\cdot\vec{v}}{\vec{v}\cdot\vec{v}} \, \vec{v}}
    \onslide<+->{\;=\; \frac{9}{5} \vetor{1}{2}}
    \onslide<+->{\;=\; \Vetor{\frac{9}{5}}{\frac{18}{5}}}
  \]
\end{resolution}

%------------------------------------------------------------------------------%
\begin{resolution}{Componente Perpendicular}
  \begin{align*}
    \onslide<+->{\vec{w}_\perp}
    \onslide<+->{& =   \vec{w} - \projuv{w}{v}}
    \\
    \onslide<+->{& =   \vetor{5}{2} - \Vetor{\frac{9}{5}}{\frac{18}{5}} }
    \onslide<+->{\;=\; \Vetor{5-\frac{9}{5}}{2-\frac{18}{5}}}
    \onslide<+->{\;=\; \Vetor{\frac{5\times 5}{5} - \frac{9}{5}}
                             {\frac{2\times 5}{5} - \frac{18}{5}}}
    \\[3mm]
    \onslide<+->{& =   \Vetor{\frac{25-9}{5}}{\frac{10-18}{5}}}
    \onslide<+->{\;=\; \Vetor{\frac{16}{5}}{\frac{-8}{5}}}
    \onslide<+->{\;=\; -\frac{8}{5}\vetor{2}{1}}
  \end{align*}
\end{resolution}

%------------------------------------------------------------------------------%
\endinput

%------------------------------------------------------------------------------%
\begin{question}
  Dados
  \[
    \vec{w} = \vetortri{2}{-1}{3}
    \qquad
    \vec{v} = \vetortri{1}{2}{2}
  \]
  Determine
  \[
    \projuv{w}{v}
  \]
\end{question}

%------------------------------------------------------------------------------%
\begin{resolution}{Solução}
\begin{columns}[t]
\column{0.5\textwidth}
  \begin{align*}
    \onslide<+->{\vec{w} \cdot \vec{v}}
    \onslide<+->{& = \vetortri{2}{-1}{3} \cdot \vetortri{1}{2}{2} \\}
    \onslide<+->{& = 2 \times 1 + (-1)2 + 3 \times 2 \\}
    \onslide<+->{& = 2 - 2 + 6}
    \onslide<+->{  = 6}
  \end{align*}
\column{0.5\textwidth}
  \begin{align*}
    \onslide<+->{\vec{v}\cdot\vec{v}}
    \onslide<+->{& = \vetortri{1}{2}{2} \cdot \vetortri{1}{2}{2} \\}
    \onslide<+->{& = 1^2 + 2^2 + 2^2 \\}
    \onslide<+->{& = 1 + 4 + 4}
    \onslide<+->{  = 9}
  \end{align*}
\end{columns}
  \[
    \onslide<+->{\projuv{w}{v}}
    \onslide<+->{\;=\; \frac{\vec{w}\cdot\vec{v}}{\vec{v}\cdot\vec{v}} \, \vec{v}}
    \onslide<+->{\;=\; \frac{6}{9} \vetortri{1}{2}{2}}
    \onslide<+->{\;=\; \frac{2}{3} \vetortri{1}{2}{2}}
    \onslide<+->{\;=\; \vetortri{\sfrac{2}{3}}{\sfrac{4}{3}}{\sfrac{4}{3}}}
  \]
\end{resolution}

%------------------------------------------------------------------------------%
\endinput

%------------------------------------------------------------------------------%
\begin{question}
  Dados
  \[
    \vec{w} = \vetortri{2}{-1}{3}
    \qquad
    \vec{v} = \vetortri{1}{2}{2}
  \]
  Determine a decomposição
  \[
    \vec{w} = \vec{w}_{\parallel} + \vec{w}_{\perp}
  \]

\end{question}

%------------------------------------------------------------------------------%
\begin{resolution}{Solução}
  Do exercício anterior temos
  \pause
  \[
    \vec{w}_{\parallel}
    \pause
    \;=\; \projuv{w}{v}
    \pause
    \;=\; \Vetortri{\frac{2}{3}}{\frac{4}{3}}{\frac{4}{3}}
  \]
\end{resolution}

%------------------------------------------------------------------------------%
\begin{resolution}{Solução}
  \begin{align*}
    \onslide<+->{\vec{w}_{\perp}}
    \onslide<+->{\;=\; \vec{w} - \vec{w}_{\parallel}}
    \onslide<+->{& =   \vetortri{2}{-1}{3} - \Vetortri{\frac{2}{3}}{\frac{4}{3}}{\frac{4}{3}}}
    \onslide<+->{\;=\; \Vetortri{2-\frac{2}{3}}{-1-\frac{4}{3}}{3-\frac{4}{3}} \\}
    \onslide<+->{& =   \Vetortri{\frac{6-2}{3}}{\frac{-3-4}{3}}{\frac{9-4}{3}}}
    \onslide<+->{\;=\; \Vetortri{\frac{4}{3}}{-\frac{7}{3}}{\frac{5}{3}}}
  \end{align*}
\end{resolution}

%------------------------------------------------------------------------------%
\endinput
%------------------------------------------------------------------------------%
\begin{question}
  Determine a projeção de $\vec{w}$ sobre $\vec{v}$
  \[
    \vec{w} = (1,\, 2,\, 3,\,  4,\, 5)
  \]
  \[
    \vec{v} = (2,\, 0,\, 1,\, -1,\, 2)
  \]
\end{question}

%------------------------------------------------------------------------------%
\begin{resolution}{Solução}
  \begin{align*}
    \onslide<+->{\vec{w}\cdot\vec{v}}
    \onslide<+->{& = (1, 2, 3,  4, 5) \cdot (2, 0, 1, -1, 2) \\}
    \onslide<+->{& = 1 \times 2 + 2 \times 0 + 3 \times 1 + 4(-1) + 5 \times 2 \\}
    \onslide<+->{& = 2 + 0 + 3 - 4 + 10}
    \onslide<+->{\;=\; 11}
  \end{align*}
  \begin{align*}
    \onslide<+->{\vec{v}\cdot\vec{v}}
    \onslide<+->{& = (2, 0, 1, -1, 2) \cdot (2, 0, 1, -1, 2) \\}
    \onslide<+->{& = 2^2 + 0^2 + 1^2 + (-1)^2 + 2^2 \\}
    \onslide<+->{& = 4 + 0 + 1 + 1 + 4}
    \onslide<+->{\;=\; 10}
  \end{align*}
\end{resolution}

%------------------------------------------------------------------------------%
\begin{resolution}{Solução}
  \begin{align*}
    \onslide<+->{\projuv{w}{v}}
    \onslide<+->{& = \frac{\vec{w}\cdot\vec{v}}{\vec{v}\cdot\vec{v}} \, \vec{v} \\[2mm]}
    \onslide<+->{& = \frac{11}{10} (2,\, 0,\, 1,\, -1,\, 2) \\[2mm]}
    \onslide<+->{& = \left(
      \frac{11}{5},\;
      0,\;
      \frac{11}{10},\;
     -\frac{11}{10},\;
      \frac{11}{5}
    \right)}
  \end{align*}
\end{resolution}

%------------------------------------------------------------------------------%
\endinput

%------------------------------------------------------------------------------%
\begin{question}
  Uma força é representada pelo vetor
  \[
    \vec{F} = (6, -2, 3)
  \]
  e atua sobre um objeto que se desloca na direção
  \[
    \vec{d} = (2, 1, 2)
  \]
  Determine as componentes da força paralela e perpendicular à direção $\vec{d}$
\end{question}

%------------------------------------------------------------------------------%
\begin{resolution}{Solução}
\begin{columns}[t]
\column{0.5\textwidth}
  \begin{align*}
    \onslide<+->{\vec{F} \cdot \vec{d}}
    \onslide<+->{& = \vetortri{6}{-2}{3} \cdot \vetortri{2}{1}{2} \\}
    \onslide<+->{& = 6 \times 2 + (-2)1 + 3 \times 2 \\}
    \onslide<+->{& = 12 - 2 + 6\\}
    \onslide<+->{& = 16}
  \end{align*}
\column{0.5\textwidth}
  \begin{align*}
    \onslide<+->{\vec{d}\cdot\vec{d}}
    \onslide<+->{& = \vetortri{2}{1}{2} \cdot \vetortri{2}{1}{2} \\}
    \onslide<+->{& = 2^2 + 1^2 + 2^2 \\}
    \onslide<+->{& = 4 + 1 + 4 \\}
    \onslide<+->{& = 9}
  \end{align*}
\end{columns}
\end{resolution}

%------------------------------------------------------------------------------%
\begin{resolution}{Solução}
  \[
    \onslide<+->{\projuv{F}{d}}
    \onslide<+->{\;=\; \frac{\vec{F}\cdot\vec{d}}{\vec{d}\cdot\vec{d}} \, \vec{d}}
    \onslide<+->{\;=\; \frac{16}{9} \vetortri{2}{1}{2}}
    \onslide<+->{\;=\; \Vetortri{\frac{16\times 2}{9}}{\frac{16}{9}}{\frac{16 \times 2}{9}}}
    \onslide<+->{\;=\; \Vetortri{\frac{32}{9}}{\frac{16}{9}}{\frac{32}{9}}}
  \]
\end{resolution}

%------------------------------------------------------------------------------%
\begin{resolution}{Solução}
  A componente paralela da força é
  \[
    \vec{F}_{\parallel}
    \;=\; \projuv{F}{d}
    \;=\; \Vetortri{\frac{32}{9}}{\frac{16}{9}}{\frac{32}{9}}
  \]
\end{resolution}

%------------------------------------------------------------------------------%
\begin{resolution}{Solução}
  A componente perpendicular é
  \[
    \vec{F}_{\perp}
    \;=\; \vec{F} - \vec{F}_{\parallel}
    \;=\; \vetortri{6}{-2}{3} - \Vetortri{\frac{32}{9}}{\frac{16}{9}}{\frac{32}{9}}
    \;=\; \Vetortri{\frac{ 6 \times 9 - 32}{9}}
                   {\frac{-2 \times 9 - 16}{9}}
                   {\frac{ 3 \times 9 - 32}{9}}
    \;=\; \Vetortri{\frac{22}{9}}{-\frac{34}{9}}{-\frac{5}{9}}
  \]
\end{resolution}

%------------------------------------------------------------------------------%
\endinput


%------------------------------------------------------------------------------%
\section{Distância}

%------------------------------------------------------------------------------%
\begin{frame}{Projeção e Distância}
  A componente perpendicular da decomposição
  \pause
  \[
    \vec{w}_{\perp}
    \pause
    \;=\; \vec{w} - \projuv{w}{v}
  \]
  \pause
  é ortogonal a $\vec{v}$

  \pause
  Sua norma é a distância entre o ponto, $P$, final de $\vec{u}$

  \pause
  e a reta, $r$, que contém o vetor $\vec{v}$

  \pause
  \[
    \dist(P, r)
    \pause
    \;=\; \norm{\vec{w}_{\perp}}
    \pause
    \;=\; \norm{\vec{w} - \projuv{w}{v}}
  \]
\end{frame}

%------------------------------------------------------------------------------%
%------------------------------------------------------------------------------%
\begin{question}
  Determine a distância entre a extremidade de
  \[
    \vec{w} = (3, 4)
  \]
  e a reta determinada por
  \[
    \vec{v} = (1, 2)
  \]
\end{question}

%------------------------------------------------------------------------------%
\begin{resolution}{Solução}
  A distância é a norma da componente perpendicular da projeção
  de $\vec{w}$ em $\vec{v}$
  \[
    \dist
    \;=\; \norm{w_\perp}
    \;=\; \norm{w - w_{\parallel}}
    \;=\; \norm{w - \projuv{w}{v}}
  \]
\end{resolution}

%------------------------------------------------------------------------------%
\begin{resolution}{Projeção}
\begin{columns}[t]
\column{0.5\textwidth}
  \begin{align*}
    \onslide<+->{\vec{w} \cdot \vec{v}}
    \onslide<+->{& = (3, 4) \cdot (1, 2) \\}
    \onslide<+->{& = 3 \times 1 + 4 \times 2 \\}
    \onslide<+->{& = 3 + 4 \\}
    \onslide<+->{& = 7}
  \end{align*}
\column{0.5\textwidth}
  \begin{align*}
    \onslide<+->{\vec{v}\cdot\vec{v}}
    \onslide<+->{& = (1, 2) \cdot (1, 2) \\}
    \onslide<+->{& = 1^2 + 2^2 \\}
    \onslide<+->{& = 1 + 4 \\}
    \onslide<+->{& = 5}
  \end{align*}
\end{columns}
  \[
    \onslide<+->{\projuv{w}{v}}
    \onslide<+->{\;=\; \frac{\vec{w}\cdot\vec{v}}{\vec{v}\cdot\vec{v}} \, \vec{v}}
    \onslide<+->{\;=\; \frac{7}{5} \vetor{1}{2}}
    \onslide<+->{\;=\; \Vetor{\frac{7}{5}}{\frac{7}{10}}}
  \]
\end{resolution}

%------------------------------------------------------------------------------%
\begin{resolution}{Componente Perpendicular}
  \begin{align*}
    w_\perp
    \;=\; w - \projuv{w}{v}
    & =\; \vetor{3}{4} - \Vetor{\frac{7}{5}}{\frac{7}{10}}
    \;=\;\Vetor{\frac{3 \times 5 - 7}{5}}{\frac{4 \times 10 - 7}{10}} \\
    & =\;\Vetor{\frac{15 - 7}{5}}{\frac{40 - 7}{10}}
    \;=\;\Vetor{\frac{8}{5}}{\frac{33}{10}}
  \end{align*}
\end{resolution}

%------------------------------------------------------------------------------%
\begin{resolution}{Distância}

    \[
      \dist
      \;=\; \norm{w_\perp}
      \;=\; \norm{w - w_{\parallel}}
      \;=\; \norm{w - \projuv{w}{v}}
    \]
\end{resolution}

%------------------------------------------------------------------------------%
\endinput



\pause

A projeção é

\[
\operatorname{proj}_{\vec{v}}\vec{w}
=
(3,0)
\]

Logo

\[
\vec{w}_{\perp}
=
(3,4)-(3,0)
\]

\[
\vec{w}_{\perp}
=
(0,4)
\]

Portanto

\[
\operatorname{dist}
=
\|\vec{w}_{\perp}\|
\]

\[
\operatorname{dist}=4
\]

\end{frame}

%------------------------------------------------------------------------------%
%\section{Lista Mínima}
%
%%------------------------------------------------------------------------------%
%\begin{frame}{Lista Mínima}
%  \begin{enumerate}
%    \item
%  \end{enumerate}
%
%  \vfill
%  Atenção: A prova é baseada no livro, não nas apresentações
%\end{frame}

%------------------------------------------------------------------------------%
\end{document}
%------------------------------------------------------------------------------%
